\section{Một từ không xuất hiện}
\label{sec:ch10-mot-tu-khong-xuat-hien}

\index{độ trung thực!vắng tên trong khảo sát}%
Chỉ số mà mục trước để lại là chỗ khảo sát tự đề xuất lối ra, và đọc lối ra ấy
cần bắt đầu từ một quan sát về chữ. Trong toàn bộ phần thân của bài, từ trang 1
tới trang 18, chữ \enquote{faithfulness} và chữ \enquote{faithful} không xuất hiện lần
nào. Trong cả tệp PDF, chuỗi ấy xuất hiện đúng một lần, ở trang 24, bên trong
nhan đề của một tài liệu tham khảo. Chữ \enquote{ground truth} cũng không xuất hiện
lần nào~\autocite{p22whichlime}.

Đó không phải là một nhận xét về từ vựng. Một khảo sát năm 2025 về toàn bộ các biến
thể của phương pháp attribution phổ biến nhất, mà hai phần ba số biến thể ấy
sinh ra để sửa quan hệ giữa lời giải thích và mô hình, viết hết 18 trang mà
không cần tới cái tên của quan hệ ấy. Đây là lần thứ hai một khảo sát viết từ
phía những người đề xuất phương pháp không giao với từ vựng của những người
phê phán chúng:
mục~\ref{sec:ch08-chi-so-nhu-mot-dung-cu} đã phải dựng lại các họ chỉ số từ
chính các bài báo đề xuất chúng, vì khảo sát XAI mà
chương~\ref{ch:khung-hop-nhat} đọc không định nghĩa một chỉ số trung thực nào để
mà chép. Hai lần ấy không phải trùng hợp. Một người đi tìm tài liệu về độ trung
thực của LIME mà tra đúng chữ độ trung thực sẽ không gặp khảo sát này, dù đây là
chỗ tập trung nhất về đúng vấn đề ấy.

Chữ mà khảo sát dùng ở chỗ ấy là \enquote{fidelity}, sáu lần trong phần thân, theo
nghĩa mà mục~\ref{sec:ch10-nam-nhom-van-de} đã tách ra. Còn ở bảng thứ ba, bảng
liệt kê các tính chất dùng để đánh giá lời giải thích, tính chất đầu tiên mang
tên \enquote{Correctness} và được mô tả là phản ánh mức chính xác mà lời giải thích
biểu diễn quá trình ra quyết định của mô hình gốc~\autocite{p22whichlime}. Đó
đúng là định nghĩa~\ref{def:ch01-do-trung-thuc}, dưới cái tên thứ ba sau
faithfulness và fidelity.

Bảng ấy có 17 tính chất, chia làm ba vùng: 8 tính chất về nội dung, 5 về cách
trình bày, 4 về người dùng; cấu trúc lấy từ một khung đánh giá XAI có sẵn của
Nauta và cộng sự, rồi bổ sung thêm những tính chất khảo sát gặp trong quá trình
rà tài liệu~\autocite{p22whichlime}. Hai tính chất trong vùng nội dung đáng chỉ
ra vì chúng là định nghĩa~\ref{def:ch10-tinh-on-dinh} bị chẻ đôi:
\enquote{Consistency} hỏi lời giải thích có ổn định qua nhiều điểm tương tự hay
không, còn \enquote{Continuity} hỏi thay đổi nhỏ ở đầu vào có dẫn tới thay đổi nhỏ
ở lời giải thích hay không.

Điều bảng ấy không có mới là chỗ nó gặp lại chương~\ref{ch:do-do-trung-thuc}.
Mỗi tính chất đi kèm một câu mô tả bằng lời và không đi kèm gì khác: không công
thức, không ngưỡng, không thủ tục tính. Vì vậy sách không phát biểu công thức
nào thay nó, đúng như chương~\ref{ch:chi-so-khong-do-do-trung-thuc} đã làm với
tám chỉ số của bài báo neo. Bảng 17 dòng ấy là một danh mục thuật ngữ, không phải
một bộ dụng cụ đo.

Với chỗ đứng ấy, đề xuất mà khảo sát đưa ra cho tương lai đọc theo đúng nghĩa
của nó. Trong bốn hướng nghiên cứu các tác giả nêu ở cuối bài, hướng thứ hai nói
rằng đã có sẵn một phổ rộng các chỉ số phù hợp và đã được thiết lập tốt, chỉ
là chúng được dùng không nhất quán trong cộng đồng, nên điều nên làm là xây một
khung đánh giá riêng cho LIME, lấy cảm hứng từ các khung đánh giá XAI đã
có~\autocite{p22whichlime}.

Sách không đọc đề xuất ấy như khảo sát đọc nó, và chỗ khác nhau phải nói rõ là
của sách chứ không phải của bài báo. Các chỉ số được gọi ở đó là đã được thiết
lập tốt chính là các họ chỉ số mà chương~\ref{ch:do-do-trung-thuc} đã liệt kê,
và điều chương ấy kết luận về chúng là chúng chưa từng được đối chiếu với một
nguồn nào ngoài chính chúng. Cụm từ \enquote{đã được thiết lập tốt} vì thế đang
nói về mức được dùng nhiều, không nói về mức đã được kiểm định. Một khung đánh
giá dựng từ những chỉ số ấy sẽ xếp hạng 48 biến thể bằng một bộ dụng cụ mà
không ai biết có đo đúng đại lượng ghi trên mặt đồng hồ hay không.

Ở đây cần nói rõ một điều để không lấy nhầm bằng chứng.
Chương~\ref{ch:chi-so-khong-do-do-trung-thuc} có một kết quả cho thấy chuyện gì
xảy ra khi một họ chỉ số lần đầu bị đem đối chiếu với ground truth dựng sẵn, và
kết quả ấy là điểm sát mức bốc thăm. Nhưng tám chỉ số bị chấm ở đó là chỉ số
dành riêng cho \tn{chuỗi suy luận}{chain-of-thought}, không phải các họ chỉ số của
chương~\ref{ch:do-do-trung-thuc}, và mục~\ref{sec:ch09-doi-tuong-khac-nhau} tồn
tại để chặn đúng phép suy rộng ấy. Cái đi qua được ranh giới là câu hỏi, không
phải con số: với các chỉ số của attribution đặc trưng, phép đối chiếu ấy chưa
được làm lần nào. Vậy nên phản bác của sách đối với đề xuất của khảo sát không
phải là các chỉ số ấy đã bị chứng minh là hỏng, mà là chưa ai kiểm chứng chúng
và khung đánh giá được đề xuất kia không có bước nào kiểm chứng chúng. Khảo sát
không nêu điều này và cũng không thể nêu: bài báo neo ra sau nó hơn một năm.
